% Chapter 1 can be included from a larger thesis with \input or \include.
% Required bibliography entries are provided in references.bib.

\chapter{Introduction}
\label{chap:introduction}

\section{Background and Motivation}
\label{sec:intro_background}

Human movement is closely coupled to sound. Musical rhythm suggests when a
movement should occur, tempo influences its speed, and timbre, dynamics, and
phrasing shape its perceived character. Audio is therefore an attractive
control modality for interactive characters and humanoid robots: it is
temporally dense, naturally available in performance settings, and more
expressive than a small set of discrete commands. A robot that responds
coherently to live music also provides a useful test case for multimodal
perception, real-time decision making, whole-body motion control, and
human--robot interaction.

Music-driven motion is nevertheless a difficult computational problem. First,
the relationship between music and dance is many-to-many, and compatibility
depends on rhythm, timbre, activity, and style rather than tempo alone. A
humanoid also imposes joint, collision, and support constraints that do not
apply to an unconstrained animated skeleton. In a live setting, these issues
must be addressed causally and within bounded latency while avoiding visible
discontinuities when the music or selected movement changes.

Research has progressed from paired music--motion benchmarks and generative
models such as AIST++/FACT and Bailando \cite{li2021ai,siyao2022bailando} to
diffusion-based choreography with improved diversity and control
\cite{alexanderson2023listen,tseng2023edge,li2023finedance,huang2024beat}.
More recent systems address limitations that are especially relevant to this
dissertation: DiscoForcing formulates dance generation as a causal,
bounded-latency streaming problem \cite{ji2026discoforcing}, while RoboPerform
generates executable humanoid actions directly and avoids an explicit
human-motion retargeting stage \cite{li2026you}. Together, these studies show
the potential of learned music-conditioned motion, but they also emphasise the
data, training, latency, and embodiment challenges involved in deploying a
real-time humanoid system.

Direct generative control is not, however, the only useful design point. It
requires substantial paired data, policy training, and a validated control
stack. Where the available actions are known in advance, retrieval offers a
complementary trade-off: motions can be inspected, retargeted, and checked for
the target robot before deployment, while selection and transition constraints
remain explicit. Although this approach cannot create new choreography, it can
provide predictable motion quality and modest runtime cost. This dissertation
therefore investigates whether a retrieval-and-control pipeline can produce
musically responsive, continuous humanoid dance from a finite, prevalidated
motion library.

\section{Research Context and Gap}
\label{sec:intro_gap}

The problem addressed in this dissertation lies at the intersection of three
research areas: cross-modal music--motion matching, real-time motion
scheduling, and humanoid motion retargeting. Each area is individually well
studied, but their integration introduces unresolved engineering and research
questions.

In cross-modal matching, a music representation must capture both semantic
similarity and movement-relevant temporal structure. A representation based
only on global style may retrieve a semantically related track while ignoring
an incompatible tempo. Conversely, a representation based only on beats per
minute may match two rhythmically similar tracks with very different energy or
timbre. A motion also has its own temporal and expressive properties, including
duration, characteristic speed, key-pose density, and the regularity of those
key poses. Consequently, retrieving a similar catalogue track is not
sufficient: the associated robot motion must also be feasible under the
speed range of the online controller and appropriate for the activity of the
query music.

In real-time scheduling, immediately selecting the highest-scoring motion at
every retrieval step can lead to rapid oscillation among near-equal candidates.
Requiring excessive evidence before switching has the opposite failure mode:
the robot continues an outdated movement after the music has changed. Even
after a target motion has been selected, entering it at an arbitrary phase can
produce joint, root, or contact discontinuities. A useful system must therefore
separate four decisions: which music is most similar, which motion is most
compatible, whether the evidence is stable enough to change motion, and where
and when the new motion should begin.

In humanoid retargeting, human and robot embodiments differ in bone lengths,
degrees of freedom, joint limits, and mechanical geometry. Retargeting errors
such as foot sliding, ground penetration, self-intersection, and abrupt
velocity spikes can degrade both perceptual fidelity and downstream control.
The study accompanying General Motion Retargeting (GMR) demonstrates that the
quality of the retargeted reference materially affects humanoid motion tracking
performance \cite{araujo2025retargeting}. This observation motivates an offline
preflight stage in which candidate actions are transformed into the Unitree G1
joint space and rejected before they can enter the online catalogue if their
metadata, joint order, continuity, ground contact, or collision properties are
invalid.

The literature therefore leaves a practical gap between offline
music-to-dance generation and a lightweight, inspectable real-time humanoid
system. Specifically, there is value in a system that:

\begin{enumerate}
    \item uses only current and past audio while maintaining a rolling musical
          context;
    \item combines semantic music similarity with explicit rhythm, tempo, and
          motion compatibility;
    \item suppresses unstable changes without eliminating musically acceptable
          diversity;
    \item aligns musical events to salient poses without instantaneously
          teleporting the motion phase;
    \item chooses transition entry states using pose, velocity, foot-contact,
          and root-motion information; and
    \item outputs target-robot joint trajectories under explicit timing,
          kinematic, and collision constraints.
\end{enumerate}

This dissertation addresses that gap through the design and evaluation of the
\emph{Humanoid Matcher}, a hybrid system comprising deterministic digital
signal processing, optional learned audio embeddings, similarity-based
retrieval, rule-governed temporal selection, and constrained kinematic motion
control.

\section{Problem Statement}
\label{sec:intro_problem}

Let a causal audio stream observed up to time $t$ be denoted by
$\mathcal{A}_{\leq t}$, and let the robot-ready motion library be

\begin{equation}
    \mathcal{D} = \{(x_i, r_i, p_i)\}_{i=1}^{N},
\end{equation}

where $x_i$ is the music descriptor associated with catalogue item $i$,
$r_i$ is a retargeted Unitree G1 trajectory, and $p_i$ contains motion metadata
such as original tempo, duration, activity, and salient key-pose phases. At
runtime, the system must select a motion index $i_t$, a phase $\phi_t$, and a
continuous robot target configuration $\mathbf{q}_t$.

The objective is not to reproduce a unique ground-truth choreography. Instead,
the system must balance five partially competing criteria:

\begin{description}
    \item[Musical relevance:] the selected motion should be associated with
    music that is similar in content, rhythm, and timbre to the live query.

    \item[Temporal compatibility:] the motion should remain usable within the
    permitted speed range, and its salient poses should be compatible with the
    observed beat and onset density.

    \item[Responsiveness and stability:] the system should react to a genuine
    musical change while avoiding oscillation caused by small ranking changes.

    \item[Motion continuity:] switching should minimise discontinuities in
    joint position, joint velocity, support state, root pose, and root velocity.

    \item[Runtime feasibility:] audio analysis, retrieval, transition planning,
    and motion playback must operate within the control-loop timing budget and
    respect joint and collision constraints.
\end{description}

These criteria cannot be reduced to a single nearest-neighbour query. They
require a hierarchy in which retrieval, compatibility scoring, temporal
stabilisation, transition planning, phase control, and output limiting are
treated as related but distinct components.

\section{Research Aim and Objectives}
\label{sec:intro_aim}

The aim of this dissertation is to design and evaluate a real-time system that
matches streaming music to a library of prevalidated humanoid dance motions and
plays the selected motions on a Unitree G1 model with beat-aware timing and
continuous transitions.

To achieve this aim, the following objectives are defined:

\begin{enumerate}
    \item construct a searchable catalogue from paired AIST++ audio and motion
          data using gain-robust audio descriptors and motion metadata;

    \item transform AIST++ SMPL sequences into a canonical Unitree G1
          representation using GMR and validate the resulting trajectories in
          MuJoCo;

    \item develop a hierarchical matcher that first retrieves related music
          and then ranks only preflight-passed motions according to tempo,
          key-pose, and activity compatibility;

    \item design a temporal selection policy that combines relevance-driven
          switching with bounded-hold diversity while avoiding ranking jitter;

    \item align accepted musical beats with salient motion poses through
          bounded speed adaptation and gradual phase correction;

    \item select and blend transition states using joint, contact, and root
          information, followed by joint-dynamics and collision checks; and

    \item evaluate retrieval accuracy, robustness, phase alignment, transition
          continuity, computational latency, and real-time scheduling behaviour.
\end{enumerate}

\section{Research Questions}
\label{sec:intro_questions}

The work is organised around the following research questions:

\begin{description}
    \item[RQ1:] To what extent can a short, gain-normalised rolling audio window
    retrieve musically related catalogue tracks and generalise to audio not
    contained in the AIST++ catalogue?

    \item[RQ2:] Does combining track similarity with tempo, motion-keypoint,
    and activity compatibility produce more suitable humanoid motion rankings
    than music similarity alone?

    \item[RQ3:] Can consecutive-evidence selection, score margins, bounded-hold
    diversity, and beat-boundary scheduling reduce unstable switching while
    remaining responsive to changes in the music?

    \item[RQ4:] Can adaptive strong-beat selection and gradual phase correction
    improve alignment between musical events and motion key poses without
    violating the configured speed range or causing phase discontinuities?

    \item[RQ5:] Can state-aware entry selection, root-aligned blending, motion
    preloading, and output limiting maintain transition continuity while
    meeting the real-time control-loop budget?
\end{description}

\section{Overview of the Proposed System}
\label{sec:intro_system}

The proposed system contains an offline preparation path and an online control
path. Offline, AIST++ motion files are paired with their synchronised music.
Audio is divided into six-second windows and encoded using a gain-robust
descriptor comprising a content embedding and explicit rhythm--timbre
features. The current catalogue contains 411 motion profiles associated with
60 music identifiers and 348 audio segments. For each motion, salient poses
are identified as prominent valleys in a whole-body velocity signal derived
from SMPL joint rotation and root translation. The human sequence is then
retargeted through GMR to the 29-degree-of-freedom Unitree G1 representation.
Only canonical trajectories that pass metadata, continuity, joint-range,
grounding, and collision-related checks are admitted to online selection.

Online, microphone input or a virtual microphone supplies both a rolling audio
window and low-level beat events. A hierarchical matcher first compares the
query with catalogue audio segments and aggregates segment evidence at the
track level. It then ranks motions linked to the best tracks using tempo,
key-pose-density, and activity compatibility. The selection policy requires
stable evidence before a relevance-driven change and can rotate among
near-optimal motions after a bounded number of held bars to preserve diversity.
Candidate motions are loaded and prepared asynchronously.

At a switch boundary, the system evaluates possible entry phases near salient
poses. The cost accounts for normalised joint-pose difference, joint-velocity
difference, foot-contact mismatch, root-state difference, and motion salience.
The selected target is aligned to the current world root and blended over a
duration derived from joint speed and acceleration limits and quantised to a
half-beat interval. During playback, accepted beats update the target phase
rate and a bounded controller gradually consumes alignment error. A final
output layer constrains joint dynamics and applies the configured runtime
collision policy before the pose is written to MuJoCo.

This architecture differs from FACT, Bailando, EDGE, RoboPerform, and
DiscoForcing in that it does not train an online motion generator. Instead, it
performs causal retrieval and control over a finite library of robot-ready
trajectories. The comparison is therefore not a claim that retrieval replaces
generative modelling. Rather, the system explores an alternative operating
point in which interpretability, predictable motion provenance, explicit
constraints, and modest runtime cost are prioritised over the ability to
synthesise novel choreography.

\section{Contributions}
\label{sec:intro_contributions}

The main contributions of this dissertation are as follows:

\begin{enumerate}
    \item A reproducible AIST++ music--motion catalogue for Unitree G1 is
          constructed, combining gain-robust audio features, motion key-pose
          statistics, GMR artefact provenance, and MuJoCo preflight results.

    \item A hierarchical, interpretable music-to-motion matching method is
          developed. It separates track-level musical similarity from
          motion-level tempo, key-pose-density, and activity compatibility, so
          that each component can be evaluated and ablated independently.

    \item A stable online selection policy is introduced that combines
          consecutive relevance evidence, a score margin relative to the
          current motion, bounded-hold diversity, least-recently-used rotation,
          beat-boundary switching, and asynchronous candidate preparation.

    \item A beat-to-key-pose controller is integrated with causal relative-beat
          contrast, adaptive weak-beat rejection, bounded speed adaptation, and
          gradual phase-error correction, avoiding instantaneous phase resets.

    \item A state-aware transition mechanism is developed that searches salient
          entry phases using pose, velocity, contact, root, and musical terms;
          aligns root trajectories; and blends motions under joint speed and
          acceleration constraints.

    \item An evaluation and diagnostic framework is provided for measuring
          retrieval recall, gain invariance, waveform-to-match latency, phase
          error, beat acceptance, transition continuity, output-limiter
          activation, collision overrides, memory usage, and real-time deadline
          misses.
\end{enumerate}

\section{Scope and Limitations}
\label{sec:intro_scope}

The dissertation focuses on real-time music matching and motion playback in a
MuJoCo model of the Unitree G1. The current player writes the floating-base pose
and named joint positions directly to the simulator state and calls forward
kinematics. It is therefore a kinematic preview and evaluation environment, not
a torque-controlled balance policy. No claim is made that the complete matcher
has been deployed on physical G1 hardware. References to robot feasibility in
this work concern joint-space validity, grounding, continuity, collision
handling, and compatibility with the G1 model rather than demonstrated dynamic
stability on hardware.

The system selects existing trajectories and does not generate new motion
content. Its expressiveness is bounded by the catalogue and by the quality of
the AIST++ source motions and GMR retargeting. The default switching policy
also assumes four accepted beats per bar; non-quadruple metre, uncertain
downbeats, ambient sound, and music with weak rhythmic structure are deliberate
stress cases rather than solved assumptions. Finally, the checked-in catalogue
uses a frozen Discogs EffNet ONNX embedding and style-tag configuration. The
local digital-signal-processing embedding remains a supported fallback, but it
must be evaluated using a separately rebuilt and frozen catalogue.

These boundaries are intentional. They allow the dissertation to isolate the
quality of music--motion retrieval, timing control, and motion continuity
without conflating those results with learned balance control or sim-to-real
policy performance.

\section{Thesis Organisation}
\label{sec:intro_organisation}

The remainder of this dissertation is organised as follows. Chapter~2 reviews
music information retrieval, music-conditioned dance synthesis, real-time
character control, humanoid motion retargeting, and motion-transition methods.
Chapter~3 defines the system requirements and presents the offline and online
architecture. Chapter~4 describes catalogue construction, audio descriptors,
motion key-pose extraction, AIST++--GMR--G1 retargeting, and trajectory
preflight. Chapter~5 presents hierarchical matching, stable motion selection,
beat--phase control, transition-entry optimisation, root continuity, and output
constraints. Chapter~6 reports experiments on retrieval, unseen-music
robustness, phase alignment, transition quality, retargeting validity, and
real-time performance. Chapter~7 concludes the dissertation, discusses the
limitations of retrieval-based kinematic playback, and outlines extensions
towards learned motion generation, metre-aware scheduling, dynamic whole-body
control, and physical humanoid deployment.
